\par\noindent\textbf{(1)} 无吸收单通道电流守恒要求 \(|S(E)|=1\)，故可写 \(S=e^{2i\delta(E)}\)，其中相移 \(\delta\) 只在模 \(\pi\) 意义下定义。若势关于原点对称，可进一步分偶宇称与奇宇称两个独立本征通道，各有自己的相移。
\par\noindent\textbf{(2)} 对窄能量波包，出射相位增加 \(2\delta(E)\)，在驻相近似下相对自由传播的时间延迟为 \(\tau_W=2\hbar\,d\delta/dE\)。同样可写 \(\tau_W=-i\hbar S^\dagger(dS/dE)\)，对单通道 \(S=e^{2i\delta}\) 后二者相同。
\par\noindent\textbf{(3)} 窄共振附近可取相移 \(\delta(E)\approx\delta_{\rm bg}+\arctan[(\Gamma/2)/(E_r-E)]\) 的连续支。对共振项求导得 \(d\delta/dE=(\Gamma/2)/[(E-E_r)^2+(\Gamma/2)^2]\)，因此峰值 \(\tau_W(E_r)\approx4\hbar/\Gamma\)（忽略背景相移斜率）。
\par\noindent\textbf{(4)} 一般势散射的 \(\tau_W\) 可以在某些频带为负，但负相对延迟并不意味着信号最前沿超因果传播；有限程势满足Wigner因果界。若存在吸收或多通道非对角散射，单相移公式不再适用，需使用完整Smith矩阵。
